\section{Compression from bounds on excitation components}
\label{sec:replica_excitation_compression}

The passage from individual excitation components to the full compressed
operator in \cite[Section~7, lines~577--590]{OpenAI2026AreaLaw} uses
Cauchy--Schwarz and the number of retained subsets. The following statements
give this implication with the component estimate as an explicit hypothesis.

\begin{theorem}[Cauchy--Schwarz for a positive quadratic form]
  \label{thm:positive_form_finite_sum}
  \lean{Matrix.PosSemidef.re_dotProduct_sum_mulVec_le_card_sum}
  \leanok
  Let $H$ be positive semidefinite on a finite-dimensional complex Hilbert
  space, let $S$ be a finite set, and let $(v_i)_{i\in S}$ be any vectors.
  Then
  \begin{align}
    \operatorname{Re}\left\langle\sum_{i\in S}v_i,
      H\sum_{i\in S}v_i\right\rangle
    \le |S|\sum_{i\in S}\operatorname{Re}\langle v_i,Hv_i\rangle.
  \end{align}
\end{theorem}
\begin{proof}
  \leanok
  Write $H=R^*R$ with $R=\sqrt H$. The left-hand side is
  $\|\sum_i Rv_i\|^2$. The triangle inequality, followed by scalar
  Cauchy--Schwarz, bounds it by $|S|\sum_i\|Rv_i\|^2$.
\end{proof}

\begin{theorem}[A bound on vectors in the defect cutoff]
  \label{thm:replica_cutoff_component_bound}
  \lean{Matrix.re_dotProduct_replicaDefectCutoff_le_of_component_bounds}
  \leanok
  \uses{def:replica_excitation_projection,def:replica_operators}
  Let $\Omega\in\mathcal H$ be a unit vector, let $k\ge0$ be an integer,
  and let $\mathcal K$ be any finite-dimensional complex Hilbert space.
  Let $N$ be the replica defect count, and let $R_B$ be the exact
  excitation projection for $B\subseteq\{0,\ldots,k-1\}$.
  Fix $0\le\tau\le1/2$, $c\ge0$, and a positive semidefinite operator
  $H$ on $\mathcal H^{\otimes k}\otimes\mathcal K$. Set
  \begin{align}
    Q_\tau&=\mathbf1_{(-\infty,\tau k]}(N)\otimes \Id_{\mathcal K},
    &u_B&=(R_B\otimes \Id_{\mathcal K})u.
  \end{align}
  Suppose $Q_\tau u=u$ and, for every $B$ with $|B|\le\tau k$,
  \begin{align}
    \operatorname{Re}\langle u_B,Hu_B\rangle\le c\|u_B\|^2.
  \end{align}
  Then, with $h(\tau)=-\tau\log\tau-(1-\tau)\log(1-\tau)$ and
  $0\log0=0$,
  \begin{align}
    \operatorname{Re}\langle u,Hu\rangle
      \le (k+1)e^{kh(\tau)}c\|u\|^2.
  \end{align}
\end{theorem}
\begin{proof}
  \leanok
  \uses{thm:positive_form_finite_sum,thm:replica_cutoff_decomposition,thm:excitation_subset_entropy}
  The actual cutoff decomposition gives $u=\sum_{|B|\le\tau k}u_B$
  and $\sum_B\|u_B\|^2=\|u\|^2$. Apply
  Theorem~\ref{thm:positive_form_finite_sum}, then the component estimate.
  The sum of the retained squared norms is at most $\|u\|^2$, and
  Theorem~\ref{thm:excitation_subset_entropy} bounds the number of retained
  subsets by $(k+1)e^{kh(\tau)}$.
\end{proof}

\begin{theorem}[Compression to a subspace of the cutoff]
  \label{thm:replica_excitation_compression_bound}
  \lean{Matrix.compression_le_of_replicaExcitationComponent_bounds}
  \leanok
  Retain the data $\Omega,k,\mathcal K,\tau,c,H,Q_\tau$ above, and let
  $P$ be an orthogonal projection satisfying $Q_\tau P=P$. For each
  vector $u$ with $Pu=u$, set $u_B=(R_B\otimes\Id_{\mathcal K})u$.
  Suppose that for every such $u$ and every $B$ with $|B|\le\tau k$,
  \begin{align}
    \operatorname{Re}\langle u_B,Hu_B\rangle\le c\|u_B\|^2.
  \end{align}
  Then
  \begin{align}
    PHP\le (k+1)e^{kh(\tau)}cP.
  \end{align}
  No commutation of $H$ with $P$, $Q_\tau$, or the individual excitation
  projections is assumed.
\end{theorem}
\begin{proof}
  \leanok
  \uses{thm:replica_cutoff_component_bound}
  For an arbitrary $x$, the vector $u=Px$ lies in the range of $P$ and
  satisfies $Q_\tau u=u$. Apply
  Theorem~\ref{thm:replica_cutoff_component_bound} to this vector.
  Orthogonality of $P$ identifies $\langle x,PHPx\rangle$ with
  $\langle Px,HPx\rangle$ and $\langle x,Px\rangle$ with $\|Px\|^2$.
  The resulting quadratic-form inequality proves the operator bound.
\end{proof}
